\begin{abstract}
How does representation geometry vary across languages, model types,
and transformer depth? We study layer-wise anisotropy in two embedding
and two generative models from the Gemma and Qwen families across
all 22 scheduled Indian languages and English. Using the aligned
IN22-Gen dataset, we sample 16,000 non-special tokens per language
and model from 1,024 sentences and evaluate IsoScore, average cosine
similarity, maximum explainable variance, and intrinsic dimensionality
at every layer. We separately examine representations before and after
the final root-mean-square normalization (RMSNorm). This distinction
substantially changes the observed geometry: in EmbeddingGemma,
final normalization raises mean token IsoScore from 0.001 to 0.228
and reduces average cosine similarity from 0.969 to 0.140. After final normalization, the two embedding
models have higher mean IsoScore than the generative models. English has the highest
group-average IsoScore in all four models, while Manipuri and Santali
form the least isotropic group in three. These findings show that
multilingual representation geometry depends on both language
and the extraction point within a model, making explicit treatment of
final normalization essential for interpreting isotropy comparisons.
\end{abstract}